\section*{अध्याय \ref{ch:b3:remote-sensing-gis} --- सुदूर संवेदन और GIS}

\begin{solution}{exo:b3:remote-sensing-gis:1}
खेत: $(0.42 - 0.06)/(0.42 + 0.06) = 0.75$, घनी हरी वनस्पति, एक फ़सल
पूरी तरह पत्तियों से भरी हुई। जुता हुआ खेत: $(0.28 - 0.21)/(0.28 + 0.21) = 0.14$, लगभग
नंगी भूमि, एक ऐसी मिट्टी जो अवरक्त की ओर थोड़ी चमकीली होती है।
\end{solution}

\begin{solution}{exo:b3:remote-sensing-gis:2}
(क)~$185000/30 \approx 6170$ पिक्सल; $185000/15 \approx 12300$ पैन्क्रोमैटिक
बैंड में। (ख)~$2^{12} = 4096$ स्तर; $2^8 = 256$। (ग)~अकेले 16 दिन;
उसी कक्षा में 8 दिन के अंतर पर चलने वाले लैंडसैट~9 के साथ, 8 दिन।
\end{solution}

\begin{solution}{exo:b3:remote-sensing-gis:3}
(क)~एक व्यतिकरण-पट्टी $\lambda/2 = \qty{27.7}{mm}$ है; छह पट्टियाँ
\qty{166}{mm} \omterm{def:b3:remote-sensing-gis:insar}{दृष्टि-रेखा} में परिवर्तन के बराबर हैं, यहाँ उपग्रह से दूर की ओर, क्योंकि
द्रोणी का अवधसान हो रहा है। (ख)~$56.6/2 = \qty{28.3}{mm}$।
\end{solution}

\begin{solution}{exo:b3:remote-sensing-gis:4}
(क)~$t = 2d/c = 2\times\num{5e5}/\num{3.00e8} = \qty{3.34}{ms}$।
(ख)~$c\,\delta t/2 = \num{3.00e8}\times\num{8e-10}/2 = \qty{0.12}{m}$; जिस सतह को बहुत-से फ़ोटॉन
नापते हैं, उसकी स्थिति कहीं बेहतर ज्ञात होती है, और सटीकता
उनकी संख्या के वर्गमूल के व्युत्क्रम के अनुसार घटती है।
\end{solution}

\begin{solution}{exo:b3:remote-sensing-gis:5}
(क)~SI इकाइयों में $L = \qty{9.0e6}{W.m^{-2}.sr^{-1}.m^{-1}}$ है। $c_1 =
2hc^2 = \qty{1.191e-16}{W.m^2.sr^{-1}}$, $c_2 = \qty{0.014388}{m.K}$ तथा
$\lambda^5 = (\num{1.09e-5})^5 = \num{1.539e-25}$, $c_1/(\lambda^5L) = 86.0$,
लेकर $T_b = c_2/(\lambda\ln 87.0) = \qty{295.6}{K}$ प्राप्त होता है। (ख)~गतिज तापमान
$0.98\,B_\lambda(T) = L$ का हल है, अर्थात् $B_\lambda(T) = L/0.98$: वही
सूत्र देता है $T = \qty{296.9}{K}$, जो \qty{1.3}{K} अधिक है।
\end{solution}

\begin{solution}{exo:b3:remote-sensing-gis:6}
(क)~$R = 693/\cos\ang{35} = \qty{846}{km}$। (ख)~वास्तविक एंटीना: $R\lambda/L =
\num{8.46e5}\times0.0555/12.3 = \qty{3.8}{km}$; संश्लेषित-द्वारक: $L/2 =
\qty{6.2}{m}$, लगभग 600 गुना महीन। (ग)~संश्लेषित-द्वारक वह
पथ-लंबाई है जिस पर कोई बिंदु किरण-पुंज में बना रहता है, $R\lambda/L =
\qty{3.8}{km}$।
\end{solution}

\begin{solution}{exo:b3:remote-sensing-gis:7}
$\partial z/\partial x = (418 - 405)/60 = 0.217$ (पूर्व घटा पश्चिम),
$\partial z/\partial y = (418 - 401)/60 = 0.283$ (उत्तर घटा दक्षिण);
$|\nabla z| = 0.357$, ढलान $\arctan 0.357 = \ang{19.6}$। भूमि
उत्तर-पूर्व की ओर चढ़ती है, इसलिए उसका मुख दक्षिण-पश्चिम की ओर है: $A =
\operatorname{atan2}(-0.217, -0.283) = \ang{217}$। हिलशेड: $\cos\ang{45}
\cos\ang{19.6} + \sin\ang{45}\sin\ang{19.6}\cos(\ang{315} - \ang{217}) =
0.666 - 0.033 = 0.63$, समतल भूमि (0.71) से थोड़ा गहरा: ढलान
उत्तर-पश्चिम के प्रकाश-स्रोत से थोड़ा दूर मुड़ती है।
\end{solution}

\begin{solution}{exo:b3:remote-sensing-gis:8}
पंक्ति योग 52, 48, 38, 22; स्तंभ योग 54, 46, 38, 22; विकर्ण 133 में से
160: समग्र यथार्थता \qty{83.1}{\percent}। उपयोगकर्ता यथार्थताएँ: जल
$48/52 = 0.92$, वन $37/48 = 0.77$, फ़सलें $29/38 = 0.76$, नगर
$19/22 = 0.86$। उत्पादक यथार्थता: जल $48/54 = 0.89$, वन $37/46 = 0.80$,
फ़सलें $29/38 = 0.76$, नगर $19/22 = 0.86$। $p_e = (52\cdot54 + 48\cdot46 +
38\cdot38 + 22\cdot22)/160^2 = 0.271$, $\kappa = (0.831 - 0.271)/(1 -
0.271) = 0.77$। फ़सलें सबसे अधिक छूटती हैं: 38 वास्तविक फ़सल-पिक्सलों में से 9
अन्यत्र गए, अधिकतर वन की ओर।
\end{solution}

\begin{solution}{exo:b3:remote-sensing-gis:9}
(क)~$L = 0.04\,B_\lambda(1443) + 0.96\,B_\lambda(290)$ का मान
\qty{10.9}{\micro\metre}, \qty{28.6}{W.m^{-2}.sr^{-1}.\micro\metre^{-1}},
पर, जिसका \omterm{def:b3:remote-sensing-gis:brightness}{दीप्ति तापमान} \qty{396}{K} है, जबकि माध्य तापमान
$0.04\times1443 + 0.96\times290 = \qty{336}{K}$ है। (ख)~
\qty{2.2}{\micro\metre}, $B_\lambda(1443)/B_\lambda(290) = \num{6.8e7}$:
इस तरंगदैर्घ्य पर लावा पिक्सल को पूरी तरह हावी करता है, और उसका \omterm{def:b3:remote-sensing-gis:brightness}{दीप्ति तापमान}
\qty{845}{K} है। लघु-तरंग अवरक्त बैंड छोटे हॉटस्पॉट को कहीं
बेहतर पहचानता है; पूरे पिक्सल के प्रति संवेदनशील ऊष्मीय बैंड वह है जो
उनके आसपास की भूमि को मापता है।
\end{solution}

\begin{solution}{exo:b3:remote-sensing-gis:10}
उपग्रह पूर्व में हो तो $d_1 = u_E\sin\theta + u_z\cos\theta$; पश्चिम में
उपग्रह हो तो $d_2 = -u_E\sin\theta + u_z\cos\theta$। जोड़ने और घटाने पर,
$u_z = (d_1 + d_2)/(2\cos\theta) = 102/1.638 = \qty{62.3}{mm}$ और $u_E =
(d_1 - d_2)/(2\sin\theta) = 26/1.147 = \qty{22.7}{mm}$। एक पास से दो अज्ञातों
के लिए एक ही समीकरण मिलता है: उत्थान और पूर्व की ओर गति का कोई भी मिश्रण, जिसका
\omterm{def:b3:remote-sensing-gis:insar}{दृष्टि-रेखा} पर प्रक्षेपण समान हो, वही व्यतिकरण-पट्टियाँ देता है।
\end{solution}

\begin{solution}{exo:b3:remote-sensing-gis:11}
(क)~\cref{thm:b3:remote-sensing-gis:fringe} से, $\Delta\phi =
4\pi\Delta r/\lambda < \pi$ के लिए आवश्यक है कि $\Delta r < \lambda/4 =
\qty{13.9}{mm}$। (ख)~$u_z = k z_0(r^2 + z_0^2)^{-3/2}$ जहाँ $k = (1 -
\nu)\Delta V/\pi$; $\mathrm du_z/\mathrm dr = -3kz_0r(r^2 + z_0^2)^{-5/2}$,
जिसका अवकलज वहाँ शून्य होता है जहाँ $(r^2 + z_0^2) = 5r^2$ है, अर्थात् $r = z_0/2$ पर। यहाँ
\[
\left|\frac{\mathrm du_z}{\mathrm dr}\right| = \frac{3k}{2z_0^3}\Big(\frac54\Big)^{-5/2}
= \frac{0.859\,k}{z_0^3} = \frac{0.859\times\num{2.39e6}}{\num{6.4e10}} = \num{3.2e-5},
\]
\qty{32}{mm} प्रति किलोमीटर। (ग)~$13.9/32 = \qty{0.43}{km}$: लगभग
\qty{430}{m} से छोटे पिक्सल, जिन्हें सेंटिनल-1 आसानी से पूरा करता है।
\end{solution}

\begin{solution}{exo:b3:remote-sensing-gis:12}
(क)~$N, R \ge 0$ के लिए, जो दोनों एक साथ शून्य न हों, $|N - R| \le N + R$, अतः $|\mathrm{NDVI}|
\le 1$; $N$ और $R$ को $aN$ और $aR$ से बदलने पर अंश और
हर को $a$ से गुणा करते हैं। (ख)~संशोधन के बिना: $(0.502 - 0.051)/(0.502 + 0.051) =
0.816$; संशोधन के साथ: $(0.512 - 0.081)/(0.512 + 0.081) = 0.727$। (ग)~पथ-विकिरण
हरी वनस्पति के सूचकांक को यहाँ लगभग एक-दसवें भाग से घटाता है, और
यह मात्रा धुंध के साथ बदलती है: असंशोधित चित्रों से बना मानचित्र
मौसम को उतना ही दिखाएगा जितना वनस्पति को, और अलग-अलग तिथियों के बीच की
तुलनाएँ निरर्थक होंगी।
\end{solution}

\begin{solution}{pb:b3:remote-sensing-gis:1}
\textbf{1.} वन $(0.38 - 0.04)/0.42 = 0.81$; 1890 का प्रवाह $0.23/0.37 =
0.62$; 1960 का $0.10/0.32 = 0.31$; 2005 का $0.02/0.26 = 0.08$।

\textbf{2.} 1890, फिर 1960, फिर 2005: सबसे पुराना प्रवाह सबसे हरा है,
आयु के क्रम में; वह अभी वन तक नहीं पहुँचा है।

\textbf{3.} नंगी चट्टान निकट अवरक्त में लाल की तुलना में थोड़ा अधिक परावर्तित करती है,
जैसे अध्याय के चित्र की भूरी रेत; लाइकेन और पहले
पौधे थोड़ा जोड़ते हैं: 0.08 लगभग नंगी भूमि का मान है।

\textbf{4.} प्रकाश से विमुख ढलान दोनों बैंडों को एक ही अनुपात में मंद कर देता है;
अनुपात इसे निरस्त कर देता है, जबकि अकेला निकट-अवरक्त परावर्तन
छायांकित वन को विरल वनस्पति समझ बैठेगा।

\textbf{5.} \cref{prop:b3:remote-sensing-gis:planck-inversion} से, तरंगदैर्घ्य
$\lambda = \qty{10.9}{\micro\metre}$ पर: पठार \qty{291.1}{K}, 
फ़ुमारोल पिक्सल \qty{302.3}{K}।

\textbf{6.} $0.97\,B_\lambda(T) = \num{8.40e6}$ से $T = \qty{293.0}{K}$ प्राप्त होता है।

\textbf{7.} $B_\lambda(\qty{370}{K}) =
\qty{22.48}{W.m^{-2}.sr^{-1}.\micro\metre^{-1}}$; $f = (9.95 - 8.40)/(22.48
- 8.40) = 0.110$, लगभग \qty{1100}{m^2}, \qty{e4}{m^2} वाले पिक्सल का।

\textbf{8.} विकिरण तापमान का एक उत्तल फलन है
(\cref{prop:b3:remote-sensing-gis:mixed}): गर्म भाग विकिरण में उसके तापमान-भाग से अधिक जोड़ता है,
अतः पिक्सल अपने माध्य से अधिक गर्म दिखता है,
$0.110\times370 + 0.890\times291.1 = \qty{299.8}{K}$ बनाम
\qty{302.3}{K}।

\textbf{9.} प्रत्येक प्रतिध्वनि की कला सटीक पथ पर निर्भर करती है; दृष्टिबिंदु में बड़ा
परिवर्तन प्रकीर्णकों की प्रतिध्वनियों के योग को बदल देता है और
\omterm{def:b3:remote-sensing-gis:insar}{व्यतिकरण-लेख} अपनी संगति खो देता है, और उच्चावच से एक कला भी जोड़ता है।
एक ही उपग्रह के $48/12 = 4$ चक्र।

\textbf{10.} $5.4\times27.73 = \qty{149.7}{mm}$ उपग्रह की ओर।

\textbf{11.} $u_z = 149.7/\cos\ang{35} = \qty{182.8}{mm}$.

\textbf{12.} भूमि ऊपर के साथ बाहर की ओर भी खिसकती है; पूर्वी ओर
बाहर की गति उपग्रह की ओर होती है और उत्थान में जुड़ती है; पश्चिमी ओर
वह घटाती है (\cref{prop:b3:remote-sensing-gis:mogi-los})। परिवर्तन
पश्चिम की ओर तेज़ी से गिरता है, जहाँ फ़्रिंज सघन हैं; प्रतिरूप
पूर्व की ओर फैला हुआ है।

\textbf{13.} उस रेखा पर क्षैतिज गति उत्तर--दक्षिण है,
दृष्टि के लंबवत है, इसलिए रडार केवल $u_z\cos\theta$ देखता है,
जो उत्थान के आकार का है। उत्थान $r = 0.766\,z_0$ पर आधा होता है, अतः $z_0 =
3.4/0.766 = \qty{4.44}{km}$।

\textbf{14.} 5.4 का आधा: 2.7 फ़्रिंज।

\textbf{15.} विरूपण दोनों व्यतिकरण-लेखों पर काल-अंतरालों से मापन के बाद
समान रहता है, या युग्म-दर-युग्म बढ़ता रहता है; वायुमंडलीय
प्रतिरूप तिथि-दर-तिथि बदलता है और प्रायः
उच्चावच का अनुसरण करता है। वही बुल्स-आई दिखाने वाला एक दूसरा युग्म
विरूपण की पुष्टि करता है।

\textbf{16.} $\Delta V = \pi u_z(0)z_0^2/(1 - \nu) = \pi\times0.1828\times
4436^2/0.75 = \qty{1.51e7}{m^3}$.

\textbf{17.} $\num{1.51e7}/48 = \qty{3.1e5}{m^3}$ प्रतिदिन,
\qty{1.1e8}{m^3} प्रति वर्ष, यदि यह दर बनी रहे।

\textbf{18.} $u_r \propto r/(r^2 + z_0^2)^{3/2}$ सबसे बड़ा वहाँ होता है जहाँ इसका
अवकलज शून्य होता है, $r^2 + z_0^2 = 3r^2$, $r = z_0/\sqrt2 = \qty{3.14}{km}$;
वहाँ $u_r = \qty{70}{mm}$, केंद्रीय उत्थान का \qty{38}{\percent}।

\textbf{19.} त्रिज्या \qty{1.48}{km}, आयतन $\frac43\pi(1.48)^3 =
\qty{13.6}{km^3}$; परिवर्तन इसका $\num{1.51e7}/\num{1.36e10} =
\qty{0.11}{\percent}$ है।

\textbf{20.} नहीं: एक फ़्रिंज का दसवाँ भाग गणना का \qty{2}{\percent} है, और
गहराई, जो शायद \qty{3.4}{km} की अर्ध-चौड़ाई से
\qty{0.1}{km} तक पढ़ी जाती है, वर्ग में प्रवेश करती है, \qty{6}{\percent}; \qty{3}{\percent}
संशोधन उस अनिश्चितता के भीतर है।

\textbf{21.} फ़ुमारोल क्षेत्र का (भाग~II) समय के साथ तापन,
उभार के नीचे छोटे भूकंपों में वृद्धि, निकली गैस में परिवर्तन,
और उपग्रह स्थिति-निर्धारण या समतलन से मापा गया उत्थान,
ढलान के कुछ बिंदुओं पर।

\textbf{22.} माउंट हैल्व्रेन के नीचे का भंडार
\textbf{15~मिलियन~\unit[text-series-to-math=true]{m^3}} ($\num{1.51e7}\,\unit{m^3}$) बढ़ गया।
\end{solution}
